\documentclass[12pt]{article}
\usepackage[utf8]{inputenc}
\usepackage[T1]{fontenc}
\usepackage[margin=1in]{geometry}
\usepackage{setspace}
\usepackage{xcolor}
\usepackage[colorlinks=true, urlcolor=blue!60!black, linkcolor=blue!60!black]{hyperref}
\setlength{\parindent}{2em}
\setlength{\parskip}{1.25em}
\renewcommand{\baselinestretch}{1.0}
\begin{document}
\begin{center}
{\LARGE \textit{NLRB v. Jones \& Laughlin Steel Corp.}, 301 U.S. 1 (1937)\par}
\vspace{0.8em}
{\large Case Brief\par}
\end{center}
\vspace{1em}

\section{Memory Jogger}

Even though manufacturing is intrastate when viewed in isolation, Congress may regulate labor relations at a vertically integrated steel plant under the Commerce Clause whenever industrial strife there would have a ``close and substantial relation to interstate commerce'' (\href{https://www.courtlistener.com/opinion/102804/national-labor-relations-board-v-jones-laughlin-steel-corp/}{CourtListener}; \href{https://supreme.justia.com/cases/federal/us/301/1/}{Justia}).

\section{Detailed Case Facts}

Congress enacted the National Labor Relations Act of July 5, 1935 (NLRA, or Wagner Act), 49 Stat. 449, 29 U.S.C. \S~151 \textit{et seq.}, to protect employees' rights to self-organization and collective bargaining (\S~7) and to prohibit employers from engaging in defined ``unfair labor practices'' (\S~8) ``affecting commerce'' (\S\S~2(6)--(7), 10(a)) (\href{https://supreme.justia.com/cases/federal/us/301/1/}{301 U.S. at 22--24, 30--33}). Section 8(1) barred interfering with, restraining, or coercing employees in exercising \S~7 rights, and \S~8(3) barred discriminating in hire or tenure of employment to encourage or discourage union membership.

Jones \& Laughlin Steel Corporation, organized under Pennsylvania law with its principal office in Pittsburgh, was the fourth-largest steel producer in the United States (\href{https://supreme.justia.com/cases/federal/us/301/1/}{301 U.S. at 25--26}; \href{https://fedsoc.org/case/nlrb-v-jones-laughlin-steel-corp}{Federalist Society}). Together with nineteen subsidiaries, it operated a completely integrated enterprise: it owned or controlled iron ore mines in Michigan and Minnesota, coal mines in Pennsylvania, limestone properties in Pennsylvania and West Virginia, four Great Lakes ore steamships, river towboats and steam barges, terminal railroads connecting its Pittsburgh and Aliquippa works to major interstate rail lines, warehouses in Chicago, Detroit, Cincinnati, and Memphis, structural steel fabricating shops in Long Island City, New York, and New Orleans, sales offices in twenty U.S. cities, and a Canadian distribution subsidiary. Approximately 75\% of its manufactured steel products were shipped out of Pennsylvania (\href{https://supreme.justia.com/cases/federal/us/301/1/}{301 U.S. at 26--27}).

At its Aliquippa, Pennsylvania plant---which employed about 10,000 workers in a town of 30,000---Jones \& Laughlin discharged ten active union leaders and officers affiliated with Beaver Valley Lodge No.\ 200 of the Amalgamated Association of Iron, Steel and Tin Workers of America, including two motor inspectors, a tractor driver, three crane operators, a coke-plant washer, and three laborers (\href{https://supreme.justia.com/cases/federal/us/301/1/}{301 U.S. at 22, 27--28}). The union filed a verified charge with the National Labor Relations Board (NLRB), which issued a complaint alleging unfair labor practices affecting commerce under \S\S~8(1) and 8(3). Jones \& Laughlin appeared specially to contest the Board's jurisdiction and constitutional authority; when its motion to dismiss was denied, the company withdrew from the hearing. The Board received evidence on the merits, found that the ten employees were discharged because of their union activity and to discourage union membership, and ordered Jones \& Laughlin to cease and desist, offer reinstatement to the ten workers, pay back wages less interim earnings, and post compliance notices for thirty days (\href{https://supreme.justia.com/cases/federal/us/301/1/}{301 U.S. at 22, 25, 29}).

\section{Procedural History}

When Jones \& Laughlin refused to comply, the NLRB petitioned the United States Circuit Court of Appeals for the Fifth Circuit (which had venue because the company operated a fabricating shop and warehouse in New Orleans) to enforce the Board's order (\href{https://supreme.justia.com/cases/federal/us/301/1/}{301 U.S. at 22, 26}). The Fifth Circuit denied the petition, 83 F.2d 998 (5th Cir. 1936), holding that the order regulated local manufacturing labor relations beyond the reach of federal power under the Commerce Clause.

The Supreme Court granted certiorari, 299 U.S. 534 (1936), and heard oral argument on February 10--11, 1937, alongside four companion ``Labor Board Cases'' (\textit{NLRB v. Fruehauf Trailer Co.}, 301 U.S. 49; \textit{NLRB v. Friedman-Harry Marks Clothing Co.}, 301 U.S. 58; \textit{Associated Press v. NLRB}, 301 U.S. 103; and \textit{Washington, Virginia \& Maryland Coach Co. v. NLRB}, 301 U.S. 142) (\href{https://supreme.justia.com/cases/federal/us/301/1/}{301 U.S. at 1, 49 n.*}). On April 12, 1937, the Supreme Court reversed the Fifth Circuit and remanded for further proceedings.

\section{Judicial Votes}

The Supreme Court reversed the Fifth Circuit by a \textbf{5--4} vote (\href{https://supreme.justia.com/cases/federal/us/301/1/}{Justia}; \href{https://ballotpedia.org/National_Labor_Relations_Board_v._Jones_\%26_Laughlin_Steel_Corporation}{Ballotpedia}):

\begin{itemize}
\item \textbf{Majority Opinion (5 Justices):} Chief Justice Charles Evans Hughes delivered the opinion of the Court, joined by Justices Louis D. Brandeis, Harlan Fiske Stone, Owen J. Roberts, and Benjamin N. Cardozo.
\item \textbf{Concurrences:} None (0).
\item \textbf{Dissenting Opinion (4 Justices):} Justice James Clark McReynolds delivered a joint dissenting opinion (published at 301 U.S. 76, covering \textit{Jones \& Laughlin}, \textit{Fruehauf Trailer}, and \textit{Friedman-Harry Marks Clothing}), joined by Justices Willis Van Devanter, George Sutherland, and Pierce Butler.
\end{itemize}

\section{Holding}

\textbf{Judgment:} Reversed and remanded. The NLRB's enforcement order against Jones \& Laughlin Steel Corporation was sustained, and the National Labor Relations Act was held constitutional as applied (\href{https://www.courtlistener.com/opinion/102804/national-labor-relations-board-v-jones-laughlin-steel-corp/}{CourtListener}; \href{https://supreme.justia.com/cases/federal/us/301/1/}{301 U.S. at 49}).

\textbf{Rule Announced:} Under the Commerce Clause, U.S. Const.\ art.\ I, \S~8, cl.\ 3, Congress may regulate intrastate activities---including manufacturing and industrial labor relations---if they have ``such a close and substantial relation to interstate commerce that their control is essential or appropriate to protect that commerce from burdens and obstructions'' (\href{https://supreme.justia.com/cases/federal/us/301/1/}{301 U.S. at 37}). In addition, the NLRA's protection of collective bargaining does not violate Fifth Amendment due process, its administrative procedures afford adequate judicial review, and its authorization of reinstatement with back pay under \S~10(c) does not violate the Seventh Amendment right to a jury trial because the statutory proceeding is not a suit at common law (\href{https://supreme.justia.com/cases/federal/us/301/1/}{301 U.S. at 43--49}).

\textbf{Holding vs.\ Dicta:} The core holding is that the NLRA validly applies to a large, nationally integrated steel manufacturer whose stoppage from labor strife would immediately and seriously disrupt interstate commerce, and that the statute's enforcement remedies survive Article III, Fifth Amendment, and Seventh Amendment challenges. Chief Justice Hughes's cautionary statement that Congress's power ``may not be extended so as to embrace effects upon interstate commerce so indirect and remote'' as to ``effectually obliterate the distinction between what is national and what is local'' (\href{https://supreme.justia.com/cases/federal/us/301/1/}{301 U.S. at 37}) was technically dicta in a decision upholding federal power---though it later became a foundational textual anchor for modern outer-limit decisions like \textit{United States v. Lopez}, 514 U.S. 549 (1995) (\href{https://scholarship.law.georgetown.edu/cgi/viewcontent.cgi?article=1433&context=facpub&sei-redir=1}{Barnett, Georgetown Law}).

\section{Analysis of Opinions}

\subsection{Majority Opinion (Chief Justice Hughes)}

Chief Justice Hughes began by invoking the canon of constitutional avoidance: ``The cardinal principle of statutory construction is to save, and not to destroy'' (\href{https://supreme.justia.com/cases/federal/us/301/1/}{301 U.S. at 30}). Rejecting Jones \& Laughlin's claim that the NLRA purported to regulate all industry nationwide, Hughes construed \S\S~2(6)--(7) and 10(a) as restricting the Board's jurisdiction strictly to unfair labor practices ``affecting'' interstate or foreign commerce, to be determined case by case.

Turning to the substantive labor rights at issue, Hughes recognized employees' right to organize and select bargaining representatives as a ``fundamental right'' comparable to a corporation's right to organize its business and choose its officers (\href{https://supreme.justia.com/cases/federal/us/301/1/}{301 U.S. at 33}). Echoing \textit{American Steel Foundries v. Tri-City Central Trades Council}, 257 U.S. 184 (1921), and the Railway Labor Act precedents (\textit{Texas \& N.O. R. Co. v. Railway Clerks}, 281 U.S. 548 (1930); \textit{Virginian Railway Co. v. System Federation No.\ 40}, 300 U.S. 515 (1937)), he explained that a single worker is dependent on daily wages and helpless against arbitrary treatment without collective organization.

On the central Commerce Clause question, Jones \& Laughlin relied on formalist precedents---including \textit{Kidd v. Pearson}, 128 U.S. 1 (1888), \textit{A.L.A. Schechter Poultry Corp. v. United States}, 295 U.S. 495 (1935), and \textit{Carter v. Carter Coal Co.}, 298 U.S. 238 (1936)---arguing that manufacturing is purely local and that raw materials stored at Aliquippa for months broke any ``stream of commerce'' (\href{https://supreme.justia.com/cases/federal/us/301/1/}{301 U.S. at 34--36}). Hughes declined to rest solely on the ``stream of commerce'' metaphor of \textit{Stafford v. Wallace}, 258 U.S. 495 (1922), calling it merely a particular illustration of Congress's plenary protective power (\href{https://digitalcommons.law.mercer.edu/cgi/viewcontent.cgi?article=2003&context=jour_mlr}{Mercer Law Review}). Instead, he announced a practical, effects-based test: ``It is the effect upon commerce, not the source of the injury, which is the criterion'' (\href{https://supreme.justia.com/cases/federal/us/301/1/}{301 U.S. at 32}). Drawing analogies to the \textit{Shreveport Rate Case}, 234 U.S. 342 (1914), and federal antitrust decisions applying the Sherman Act to manufacturers and mineworkers, Hughes held that a stoppage at Jones \& Laughlin's far-flung steel enterprise---referencing the devastating 1919--1920 steel strike---would have an ``immediate'' and potentially ``catastrophic'' impact on interstate commerce (\href{https://supreme.justia.com/cases/federal/us/301/1/}{301 U.S. at 41--43}). He distinguished \textit{Schechter} as involving remote local effects and \textit{Carter Coal} as turning on improper delegation and due process flaws.

Finally, Hughes rejected Jones \& Laughlin's Fifth Amendment, administrative-procedure, and Seventh Amendment objections. The NLRA did not compel employers to reach collective agreements or bar legitimate discharges unrelated to union coercion; its notice, hearing, substantial-evidence, and judicial-review provisions satisfied due process; and reinstatement with back pay under \S~10(c) was a statutory, equitable-type enforcement remedy rather than a common-law suit triggering the Seventh Amendment jury right (\href{https://supreme.justia.com/cases/federal/us/301/1/}{301 U.S. at 44--49}).

\subsection{Dissenting Opinion (Justice McReynolds, joined by Justices Van Devanter, Sutherland, and Butler)}

Justice McReynolds, writing for the four conservative dissenters (the ``Four Horsemen'') in a combined dissent covering \textit{Jones \& Laughlin} and two smaller manufacturing companion cases, argued that the majority had effectively abandoned the Constitution's distinction between interstate commerce and local production (\href{https://supreme.justia.com/cases/federal/us/301/58/}{301 U.S. at 76}; \href{https://fedsoc.org/case/nlrb-v-jones-laughlin-steel-corp}{Federalist Society}). Relying heavily on \textit{Schechter Poultry} and \textit{Carter Coal}, McReynolds maintained that manufacturing is distinct from commerce and remains subject to state police power.

In McReynolds's view, discharging ten production workers at a Pennsylvania plant had an effect on interstate commerce that was ``indirect and remote in the highest degree'' (\href{https://historynet.com/1937-labor-board-case-upheld-workers-right-to-organize/}{HistoryNet}; \href{https://supreme.justia.com/cases/federal/us/301/58/}{301 U.S. 58}). He warned that upholding the NLRA on the theory that a potential strike might reduce interstate shipments adopted a ``view of congressional power that would extend it into almost every field of human industry,'' destroying the federal balance and infringing employers' Fifth Amendment liberty of contract (\href{https://historynet.com/1937-labor-board-case-upheld-workers-right-to-organize/}{HistoryNet}; \href{https://supreme.justia.com/cases/federal/us/301/1/}{Justia}).

\section{Examples: Future Applications}

\begin{enumerate}
\item \textbf{Same-Side (Valid Federal Regulation):} Congress applies federal collective-bargaining protections and workplace safety standards to a mid-sized Virginia garment manufacturer that imports most of its fabric from out-of-state suppliers and ships over 80\% of its finished suits to interstate buyers. Following \textit{Jones \& Laughlin} (and its companion decision in \textit{NLRB v. Friedman-Harry Marks Clothing Co.}, 301 U.S. 58 (1937)), this falls within the Commerce Clause because industrial disruption in an enterprise tightly integrated into interstate supply and distribution markets substantially burdens interstate commerce (\href{https://supreme.justia.com/cases/federal/us/301/58/}{Justia}).

\item \textbf{Same-Side (Valid Federal Regulation):} Congress regulates labor disputes at an intrastate electric utility that operates entirely within New York but supplies power to interstate railroads, steamship terminals, and telegraph facilities. Under \textit{Jones \& Laughlin}'s practical effects test (later applied in \textit{Consolidated Edison Co. v. NLRB}, 305 U.S. 197 (1938)), federal power reaches the utility because a strike cutting off power would immediately paralyze instrumentalities of interstate commerce (\href{https://supreme.justia.com/cases/federal/us/305/197/}{Justia}).

\item \textbf{Opposite-Side (Exceeds Commerce Power):} Congress enacts a federal criminal statute prohibiting the simple possession of a firearm within 1,000 feet of any local public school, containing no jurisdictional hook to interstate commerce and regulating non-economic conduct. Under Chief Justice Hughes's explicit federalism boundary in \textit{Jones \& Laughlin}---that Congress may not reach effects ``so indirect and remote'' as to ``obliterate the distinction between what is national and what is local'' (301 U.S. at 37)---this statute exceeds the Commerce Clause, as the Court later held in \textit{United States v. Lopez}, 514 U.S. 549 (1995) (\href{https://scholarship.law.georgetown.edu/cgi/viewcontent.cgi?article=1433&context=facpub&sei-redir=1}{Barnett, Georgetown Law}).

\item \textbf{Opposite-Side (Exceeds Commerce Power):} Congress passes a statute prescribing uniform state-court standards for non-economic domestic-relations custody disputes, reasoning that family instability reduces workforce productivity and indirectly lowers national GDP. Under \textit{Jones \& Laughlin}'s warning against collapsing the dual system of government into ``a completely centralized government'' (\href{https://supreme.justia.com/cases/federal/us/301/1/}{301 U.S. at 37}), the causal chain to interstate commerce is far too attenuated to sustain federal authority.

\item \textbf{Fence-Sitter (Borderline):} Congress applies federal labor or agricultural quotas to a small, single-location family farm or neighborhood bakery that purchases minor equipment from out of state but sells all of its output locally to walk-in customers. Viewed strictly through \textit{Jones \& Laughlin}'s focus on a single enterprise's ``immediate'' and ``catastrophic'' impact on interstate commerce (301 U.S. at 41), a strike or stoppage at one small local business would appear remote and local; yet under the post-1937 aggregation expansion that \textit{Jones \& Laughlin} set in motion (culminating in \textit{Wickard v. Filburn}, 317 U.S. 111 (1942)), courts could uphold regulation if the aggregate nationwide impact of similar small producers substantially affects interstate markets (\href{https://ballotpedia.org/National_Labor_Relations_Board_v._Jones_\%26_Laughlin_Steel_Corporation}{Ballotpedia}).
\end{enumerate}

\section{Critique}

\textit{NLRB v. Jones \& Laughlin Steel Corp.} is universally recognized as a watershed moment in American constitutional history, marking the collapse of \textit{Lochner}-era formalist Commerce Clause limits and the birth of modern federal regulatory authority (\href{https://historynet.com/1937-labor-board-case-upheld-workers-right-to-organize/}{HistoryNet}).

\textbf{Originalist and Textualist Perspectives:} Originalist scholars sharply criticize \textit{Jones \& Laughlin} for untethering the Commerce Clause from its founding-era meaning. Richard Epstein argues that ``commerce'' originally meant interstate trade and exchange rather than local manufacturing or agriculture, and that replacing categorical limits with an open-ended ``substantial effects'' test effectively erased the doctrine of enumerated powers (\href{https://chicagounbound.uchicago.edu/cgi/viewcontent.cgi?referer=&httpsredir=1&article=2284&context=journal_articles}{Epstein, \textit{The Mistakes of 1937}, 11 Geo.\ Mason U.\ L.\ Rev.\ 5 (1988)}; \href{https://chicagounbound.uchicago.edu/cgi/viewcontent.cgi?article=2315&context=journal_articles}{Epstein, 73 Va.\ L.\ Rev.\ 1387 (1987)}). Similarly, Randy Barnett and Heritage Foundation commentators contend that although Hughes promised to preserve the line between national and local concerns, his pragmatic test laid the groundwork for \textit{United States v. Darby} (1941) and \textit{Wickard v. Filburn} (1942), giving Congress near-plenary police power until \textit{United States v. Lopez} (1995) (\href{https://scholarship.law.georgetown.edu/cgi/viewcontent.cgi?article=1433&context=facpub&sei-redir=1}{Barnett, Georgetown Law}; \href{https://constitution.heritage.org/essays/a1-s8-c3-b/}{Heritage Guide}). By contrast, some market-oriented originalists, such as Robert Pushaw, argue that the founding generation understood ``commerce'' more broadly than mere trade, encompassing market-oriented production and labor, making \textit{Jones \& Laughlin} more defensible on originalist grounds than critics claim (\href{https://digitalcommons.law.byu.edu/cgi/viewcontent.cgi?article=3408&context=lawreview}{Pushaw, BYU L.\ Rev.\ (2022)}).

\textbf{Progressive, Historical, and Doctrinal Perspectives:} Progressive and doctrinal scholars praise Chief Justice Hughes for abandoning sterile formalisms---such as the rigid ``direct vs.\ indirect'' and ``production vs.\ commerce'' dichotomies---in favor of economic realism suited to an integrated national industrial economy (\href{https://fedsoc.org/case/nlrb-v-jones-laughlin-steel-corp}{Federalist Society}; \href{https://digitalcommons.law.mercer.edu/cgi/viewcontent.cgi?article=2003&context=jour_mlr}{Mercer Law Review}). Historians vigorously debate the decision's political context: traditional accounts (and Bruce Ackerman's theory of ``constitutional moments'') view Justice Roberts's and Chief Justice Hughes's votes in 1937 as ``the switch in time that saved nine'' in response to President Franklin D. Roosevelt's Court-packing plan, whereas revisionist legal historians such as Barry Cushman argue that the doctrinal shift evolved organically from internal tensions in earlier Commerce Clause and due-process case law (\href{https://www.uclalawreview.org/wp-content/uploads/2019/09/48_54UCLALRev1847August2007.pdf}{UCLA Law Review (2007)}; \href{https://historynet.com/1937-labor-board-case-upheld-workers-right-to-organize/}{HistoryNet}).

\textbf{Competing Underlying Values:} The decision encapsulates the fundamental tension between \textit{national democratic capacity} (empowering Congress to address nationwide economic crises, industrial conflict, and interstate externalities in an interconnected market) and \textit{enumerated-powers federalism} (preserving meaningful constitutional boundaries on central authority and safeguarding state autonomy and individual economic liberty).

\section{Key Quotations}

\begin{quote}
``The authority of the federal government may not be pushed to such an extreme as to destroy the distinction, which the commerce clause itself establishes, between commerce `among the several States' and the internal concerns of a State. That distinction between what is national and what is local in the activities of commerce is vital to the maintenance of our federal system.''
\end{quote}
\noindent--- Chief Justice Hughes, Opinion of the Court, 301 U.S. at 30 (\href{https://supreme.justia.com/cases/federal/us/301/1/}{Justia}; \href{https://www.courtlistener.com/opinion/102804/national-labor-relations-board-v-jones-laughlin-steel-corp/}{CourtListener}).

\begin{quote}
``Although activities may be intrastate in character when separately considered, if they have such a close and substantial relation to interstate commerce that their control is essential or appropriate to protect that commerce from burdens and obstructions, Congress cannot be denied the power to exercise that control.''
\end{quote}
\noindent--- Chief Justice Hughes, Opinion of the Court, 301 U.S. at 37 (\href{https://supreme.justia.com/cases/federal/us/301/1/}{Justia}; \href{https://ballotpedia.org/National_Labor_Relations_Board_v._Jones_\%26_Laughlin_Steel_Corporation}{Ballotpedia}).

\begin{quote}
``In view of respondent's far-flung activities, it is idle to say that the effect would be indirect or remote. It is obvious that it would be immediate, and might be catastrophic. We are asked to shut our eyes to the plainest facts of our national life, and to deal with the question of direct and indirect effects in an intellectual vacuum.''
\end{quote}
\noindent--- Chief Justice Hughes, Opinion of the Court, 301 U.S. at 41 (\href{https://supreme.justia.com/cases/federal/us/301/1/}{Justia}; \href{https://www.courtlistener.com/opinion/102804/national-labor-relations-board-v-jones-laughlin-steel-corp/}{CourtListener}).

\begin{quote}
``We have often said that interstate commerce itself is a practical conception. It is equally true that interferences with that commerce must be appraised by a judgment that does not ignore actual experience.''
\end{quote}
\noindent--- Chief Justice Hughes, Opinion of the Court, 301 U.S. at 41--42 (\href{https://supreme.justia.com/cases/federal/us/301/1/}{Justia}).

\begin{quote}
``The things inhibited by the Labor Act relate to the management of a manufacturing plant---something distinct from commerce and subject to the authority of the state. And this may not be abridged because of some vague possibility of distant interference with commerce.''
\end{quote}
\noindent--- Justice McReynolds, dissenting, 301 U.S. at 99 (\href{https://supreme.justia.com/cases/federal/us/301/58/}{Justia}; \href{https://www.courtlistener.com/opinion/102804/national-labor-relations-board-v-jones-laughlin-steel-corp/}{CourtListener}).

\end{document}
